% GENERATED FILE. Edit canonical lessons, then run npm run generate:books.
% canonical-source-hash: fnv1a64:7adbf4e8aedabb1f
% canonical-lessons: TE-C352-vembadincu, TE-C352-dakko, TE-C352-longipovu, TE-C352-pattincu, TE-C352-tappuko

\chapter{The Chase}
\label{ch:te-the-chase}
\hlchaptermodality{352}

\begin{chapteropening}
\textbf{By the end of this chapter:} \emph{I can say to pursue, to hide oneself, to surrender, to turn in and to step aside.}

Complete the last lesson of chapter 352: I can say to pursue, to hide oneself, to surrender, to turn in and to step aside.
\end{chapteropening}

\section[\texorpdfstring{\te{వెంబడించు} (vembaḍiñcu)}{వెంబడించు (vembaḍiñcu)}]{\te{వెంబడించు} (vembaḍiñcu) — to pursue; to chase}

\label{lesson:TE-C352-vembadincu}



\begin{quote}
Before the new one: say the Telugu for to release, then the Telugu for to free from captivity.
\end{quote}

\subsection*{You'll want to know: \te{వెంబడించు}}
\textbf{\te{వెంబడించు}} — \emph{vembaḍiñcu} — \textquotedblleft{}to pursue; to chase\textquotedblright{}.

\begin{quote}
\te{పోలీసులు} \te{దొంగను} \te{కారులో} \te{వెంబడించారు}.
\end{quote}

\begin{quote}
\emph{pōlīsulu doṅganu kārulō vembaḍiñcāru.}
\end{quote}

\begin{quote}
The police chased the thief in a car.
\end{quote}

Past (I …): \te{వెంబడించాను}, \emph{vembaḍiñcānu}.

\te{వెంబడి} means close behind; \te{తరుము} is the rougher, more spoken word.

\begin{cousinweb}
\textbf{\te{వెంబడించు}} (\emph{vembaḍiñcu}) is native Dravidian: DEDR 5488 lists Telugu \emph{vembaḍi}, \textquotedblleft{}along with\textquotedblright{}, and \emph{vembaḍincu}, under the word for \textquotedblleft{}the back\textquotedblright{}.
\end{cousinweb}

\begin{grammarlens}[title={What you've built}]
One more word for telling what happened.
\end{grammarlens}

\subsection*{Your turn}
\begin{itemize}
\raggedright
  \item \emph{Say it:} \emph{vembaḍiñcu}
  \item \emph{Say it:} \emph{vembaḍiñcu}, once more
  \item \emph{Say it:} \emph{pōlīsulu doṅganu kārulō vembaḍiñcāru.} — The police chased the thief in a car.
  \item \emph{Recall:} say the Telugu for to release, then the Telugu for to free from captivity, then say \emph{vembaḍiñcu} again
\end{itemize}

\begin{tcolorbox}[breakable,colback=teal!4,colframe=teal!35!black,title={Before you move on}]
What does \te{వెంబడించు} mean? (\textquotedblleft{}To pursue; to chase\textquotedblright{}.) Say it once more.
\end{tcolorbox}
\section[\texorpdfstring{\te{దాక్కో} (dākkō)}{దాక్కో (dākkō)}]{\te{దాక్కో} (dākkō) — to hide oneself}

\label{lesson:TE-C352-dakko}



\begin{quote}
Before the new one: say the Telugu for to free from captivity, then the Telugu for to pursue.
\end{quote}

\subsection*{You'll want to know: \te{దాక్కో}}
\textbf{\te{దాక్కో}} — \emph{dākkō} — \textquotedblleft{}to hide oneself\textquotedblright{}.

\begin{quote}
\te{దొంగ} \te{చెట్టు} \te{వెనక} \te{దాక్కున్నాడు}.
\end{quote}

\begin{quote}
\emph{doṅga ceṭṭu venaka dākkunnāḍu.}
\end{quote}

\begin{quote}
The thief hid behind a tree.
\end{quote}

Past (I …): \te{దాక్కున్నాను}, \emph{dākkunnānu}.

\te{దాచు} is to hide a thing; \te{దాక్కో} is to hide yourself.

\begin{cousinweb}
\textbf{\te{దాక్కో}} (\emph{dākkō}) is native Dravidian: DEDR 63 lists Telugu \emph{dāgu}, \textquotedblleft{}to lie hidden\textquotedblright{}, beside \emph{dācu}, \textquotedblleft{}to hide, conceal\textquotedblright{}; \te{దాక్కో} adds the self-ending \emph{-kō}.
\end{cousinweb}

\begin{grammarlens}[title={What you've built}]
One more word for telling what happened.
\end{grammarlens}

\subsection*{Your turn}
\begin{itemize}
\raggedright
  \item \emph{Say it:} \emph{dākkō}
  \item \emph{Say it:} \emph{dākkō}, once more
  \item \emph{Say it:} \emph{doṅga ceṭṭu venaka dākkunnāḍu.} — The thief hid behind a tree.
  \item \emph{Recall:} say the Telugu for to free from captivity, then the Telugu for to pursue, then say \emph{dākkō} again
\end{itemize}

\begin{tcolorbox}[breakable,colback=teal!4,colframe=teal!35!black,title={Before you move on}]
What does \te{దాక్కో} mean? (\textquotedblleft{}To hide oneself\textquotedblright{}.) Say it once more.
\end{tcolorbox}
\section[\texorpdfstring{\te{లొంగిపోవు} (loṅgipōvu)}{లొంగిపోవు (loṅgipōvu)}]{\te{లొంగిపోవు} (loṅgipōvu) — to surrender; to give in}

\label{lesson:TE-C352-longipovu}



\begin{quote}
Before the new one: say the Telugu for to pursue, then the Telugu for to hide oneself.
\end{quote}

\subsection*{You'll want to know: \te{లొంగిపోవు}}
\textbf{\te{లొంగిపోవు}} — \emph{loṅgipōvu} — \textquotedblleft{}to surrender; to give in\textquotedblright{}.

\begin{quote}
\te{చివరికి} \te{దొంగ} \te{పోలీసులకు} \te{లొంగిపోయాడు}.
\end{quote}

\begin{quote}
\emph{civariki doṅga pōlīsulaku loṅgipōyāḍu.}
\end{quote}

\begin{quote}
In the end the thief surrendered to the police.
\end{quote}

Past (I …): \te{లొంగిపోయాను}, \emph{loṅgipōyānu}.

To whom takes -\te{కు}: \te{పోలీసులకు} \te{లొంగిపోవు}.

\begin{cousinweb}
\textbf{\te{లొంగిపోవు}} (\emph{loṅgipōvu}) is native Dravidian: DEDR 698 lists Telugu \emph{loṅgu}, \textquotedblleft{}to yield, submit\textquotedblright{}, with a query, and \emph{pōvu}, \textquotedblleft{}to go\textquotedblright{} (DEDR 4572), completes the action.
\end{cousinweb}

\begin{grammarlens}[title={What you've built}]
One more word for telling what happened.
\end{grammarlens}

\subsection*{Your turn}
\begin{itemize}
\raggedright
  \item \emph{Say it:} \emph{loṅgipōvu}
  \item \emph{Say it:} \emph{loṅgipōvu}, once more
  \item \emph{Say it:} \emph{civariki doṅga pōlīsulaku loṅgipōyāḍu.} — In the end the thief surrendered to the police.
  \item \emph{Recall:} say the Telugu for to pursue, then the Telugu for to hide oneself, then say \emph{loṅgipōvu} again
\end{itemize}

\begin{tcolorbox}[breakable,colback=teal!4,colframe=teal!35!black,title={Before you move on}]
What does \te{లొంగిపోవు} mean? (\textquotedblleft{}To surrender; to give in\textquotedblright{}.) Say it once more.
\end{tcolorbox}
\section[\texorpdfstring{\te{పట్టించు} (paṭṭiñcu)}{పట్టించు (paṭṭiñcu)}]{\te{పట్టించు} (paṭṭiñcu) — to turn in (to the police)}

\label{lesson:TE-C352-pattincu}



\begin{quote}
Before the new one: say the Telugu for to hide oneself, then the Telugu for to surrender.
\end{quote}

\subsection*{You'll want to know: \te{పట్టించు}}
\textbf{\te{పట్టించు}} — \emph{paṭṭiñcu} — \textquotedblleft{}to turn in (to the police)\textquotedblright{}.

\begin{quote}
\te{గ్రామస్తులు} \te{దొంగను} \te{పోలీసులకు} \te{పట్టించారు}.
\end{quote}

\begin{quote}
\emph{grāmastulu doṅganu pōlīsulaku paṭṭiñcāru.}
\end{quote}

\begin{quote}
The villagers turned the thief in to the police.
\end{quote}

Past (I …): \te{పట్టించాను}, \emph{paṭṭiñcānu}.

The causative of \te{పట్టు}, to catch: to get someone caught. Not the same as \te{పట్టించుకో}, to pay attention.

\begin{cousinweb}
\textbf{\te{పట్టించు}} (\emph{paṭṭiñcu}) is native Dravidian, the causative of \emph{paṭṭu}, \textquotedblleft{}to hold, catch\textquotedblright{} (DEDR 4034).
\end{cousinweb}

\begin{grammarlens}[title={What you've built}]
One more word for telling what happened.
\end{grammarlens}

\subsection*{Your turn}
\begin{itemize}
\raggedright
  \item \emph{Say it:} \emph{paṭṭiñcu}
  \item \emph{Say it:} \emph{paṭṭiñcu}, once more
  \item \emph{Say it:} \emph{grāmastulu doṅganu pōlīsulaku paṭṭiñcāru.} — The villagers turned the thief in to the police.
  \item \emph{Recall:} say the Telugu for to hide oneself, then the Telugu for to surrender, then say \emph{paṭṭiñcu} again
\end{itemize}

\begin{tcolorbox}[breakable,colback=teal!4,colframe=teal!35!black,title={Before you move on}]
What does \te{పట్టించు} mean? (\textquotedblleft{}To turn in (to the police)\textquotedblright{}.) Say it once more.
\end{tcolorbox}
\section[\texorpdfstring{\te{తప్పుకో} (tappukō)}{తప్పుకో (tappukō)}]{\te{తప్పుకో} (tappukō) — to step aside; to withdraw}

\label{lesson:TE-C352-tappuko}



\begin{quote}
Before the new one: say the Telugu for to surrender, then the Telugu for to turn in.
\end{quote}

\subsection*{You'll want to know: \te{తప్పుకో}}
\textbf{\te{తప్పుకో}} — \emph{tappukō} — \textquotedblleft{}to step aside; to withdraw\textquotedblright{}.

\begin{quote}
\te{పోలీసు} \te{జీపు} \te{రాగానే} \te{జనం} \te{తప్పుకున్నారు}.
\end{quote}

\begin{quote}
\emph{pōlīsu jīpu rāgānē janaṁ tappukunnāru.}
\end{quote}

\begin{quote}
As soon as the police jeep came, people stepped aside.
\end{quote}

Past (I …): \te{తప్పుకున్నాను}, \emph{tappukunnānu}.

Also to pull out: \te{పోటీ} \te{నుంచి} \te{తప్పుకో}, to withdraw from a contest.

\begin{cousinweb}
\textbf{\te{తప్పుకో}} (\emph{tappukō}) is native Dravidian: \emph{tappu}, \textquotedblleft{}to miss, slip, escape\textquotedblright{} (DEDR 3071), with the self-ending \emph{-kō}: to slip out of the way.
\end{cousinweb}

\begin{grammarlens}[title={What you've built}]
One more word for telling what happened.
\end{grammarlens}

\subsection*{Your turn}
\begin{itemize}
\raggedright
  \item \emph{Say it:} \emph{tappukō}
  \item \emph{Say it:} \emph{tappukō}, once more
  \item \emph{Say it:} \emph{pōlīsu jīpu rāgānē janaṁ tappukunnāru.} — As soon as the police jeep came, people stepped aside.
  \item \emph{Recall:} say the Telugu for to surrender, then the Telugu for to turn in, then say \emph{tappukō} again
\end{itemize}

\begin{tcolorbox}[breakable,colback=teal!4,colframe=teal!35!black,title={Before you move on}]
What does \te{తప్పుకో} mean? (\textquotedblleft{}To step aside; to withdraw\textquotedblright{}.) Say it once more.
\end{tcolorbox}
